\documentclass[11pt, a4paper]{article}

% XeLaTeX for Korean + Greek support
\usepackage{fontspec}
\usepackage{kotex}
\usepackage{geometry}
\usepackage{titlesec}
\usepackage{enumitem}
\usepackage{hyperref}
\usepackage{xcolor}
\usepackage{fancyhdr}

\geometry{margin=2.5cm}

% Font settings - Overleaf has Noto fonts
\setmainfont{Noto Serif}
\setsansfont{Noto Sans}
\setmonofont{Noto Sans Mono}
\setmainhangulfont{Noto Serif CJK KR}
\setsanshangulfont{Noto Sans CJK KR}

% Section formatting
\titleformat{\section}{\Large\bfseries}{}{0em}{}
\titleformat{\subsection}{\large\bfseries}{}{0em}{}
\titleformat{\subsubsection}{\normalsize\bfseries}{}{0em}{}

% Note command for translator notes
\newcommand{\tnote}[1]{\par\noindent\textit{\textcolor{gray}{[역주] #1}}\par\smallskip}

% Header/footer
\pagestyle{fancy}
\fancyhf{}
\fancyhead[L]{\small 소크라테스의 변론}
\fancyhead[R]{\small 한국어 번역}
\fancyfoot[C]{\thepage}
\renewcommand{\headrulewidth}{0.4pt}

\hypersetup{
    colorlinks=true,
    linkcolor=black,
    urlcolor=blue!60!black
}

\begin{document}

% Title page
\begin{titlepage}
\centering
\vspace*{5cm}
{\Huge\bfseries 소크라테스의 변론\par}
\vspace{0.5cm}
{\Large\textit{Ἀπολογία Σωκράτους}\par}
\vspace{2cm}
{\Large 플라톤 저\par}
\vspace{0.5cm}
{\normalsize Benjamin Jowett 영역 기반 한국어 번역\par}
\vspace{1cm}
{\small AI 번역: Claude Opus 4.6 (Anthropic)\par}
{\small 기획, 용어 선정, 검토: Junhee Hyeon\par}
\vspace{0.5cm}
{\small Jowett 영역본 기반, 핵심 용어 그리스어 원어 참조\par}
\vfill
{\small 2026년 3월\par}
\end{titlepage}

\tableofcontents
\newpage

\section*{번역자 주}


이 번역은 Benjamin Jowett의 영역본을 기반으로 합니다. Jowett 번역은 고전적이지만 빅토리아 시대 영어의 수사적 장식이 가미되어 있어, 고대 그리스어 원문의 뉘앙스와 반드시 일치하지는 않습니다. 특히 'God'이라는 단어는 원문에서 'ho theos(신)' 또는 'to theion(신적인 것)'으로 다양하게 표현되는데, Jowett은 이를 단일한 'God'으로 통일하는 경향이 있다는 지적이 있습니다(이것이 의도적 기독교화인지 19세기 영어 관습의 반영인지는 학자마다 견해가 다릅니다). 이런 부분은 역주로 표시했습니다.


'Apology(변론)'는 사과가 아니라 '변호/변론'을 뜻하는 그리스어 'apologia(ἀπολογία)'에서 온 것입니다. 이 번역은 고대 그리스어 원문에서 직접 번역한 것이 아니라, Benjamin Jowett의 영어 번역본을 기반으로 AI(Claude Opus 4.6, Anthropic)가 한국어로 옮긴 것입니다. 따라서 원문 → 영역 → 한국어의 이중 번역 과정에서 의미가 변형되었을 수 있습니다. 역주 역시 AI의 지식에 기반한 해석이며, 전문적인 고전학 연구를 대체하지 않습니다.


\bigskip


\newpage
\section{제1부: 소크라테스의 변론}


\subsection{서두 — 자신의 말하기 방식에 대하여}


아테네 시민 여러분, 고발자들의 연설을 듣고 여러분이 어떤 느낌을 받았는지 나는 알 수 없습니다. 다만 그들의 설득력 있는 말이 나마저도 내가 누구인지 잊게 만들 뻔했다는 것은 압니다 — 그만큼 효과적이었습니다. 하지만 그들은 진실에 해당하는 말을 거의 한마디도 하지 않았습니다.


그들의 수많은 거짓 중 특히 놀라운 것이 하나 있었는데, 나의 언변(eloquence)에 속지 말라고 여러분에게 경고한 것입니다. 그들이 이 말을 부끄러워해야 하는 이유는, 내가 입을 여는 순간 나의 부족함이 드러나 그 말이 거짓임이 바로 밝혀질 것이기 때문입니다. 다만 그들이 말하는 '언변의 힘'이 '진실의 힘'을 뜻한다면, 그렇다면 나도 과연 언변이 있다고 인정하겠습니다. 하지만 그들과는 전혀 다른 종류의 것입니다!


그들은 진실을 거의 한마디도 말하지 않았지만, 여러분은 나에게서 온전한 진실(the whole truth)을 듣게 될 것입니다. 다만 그들처럼 꾸며진 연설문으로 전하지는 않을 것입니다. 아닙니다! 나는 그 순간에 떠오르는 말과 논변을 사용할 것입니다. 이것이 옳다고 확신하기 때문입니다. 나이 칠십이 넘어 처음으로 법정에 선 사람에게 젊은 연설가 같은 것을 기대하지 마십시오.


한 가지 부탁이 있습니다. 내가 변론에서 평소 쓰던 말을 사용하더라도 — 여러분 대부분이 아고라에서, 환전상의 탁자 곁에서, 또는 다른 곳에서 들었을 바로 그 말을 — 놀라거나 방해하지 말아주십시오. 나는 70세가 넘었고, 이것은 내 생애 처음으로 법정에 선 일입니다. 나를 타국에서 온 이방인이라 생각하고, 그가 자기 모국어로, 자기 나라의 방식대로 말하는 것을 양해해 주듯이 해주십시오.


\textbf{말하는 방식은 신경 쓰지 마시고, 내 주장의 정의로움만을 생각하십시오. 재판관은 정의롭게 판단하고, 변론자는 진실을 말하는 것이 각자의 덕(virtue)입니다.}


\tnote{'virtue'는 그리스어 aretē(ἀρετή)의 번역인데, 이는 '탁월함/본래의 기능을 잘 수행함'에 가깝지 도덕적 '덕'만을 의미하지는 않음.}


\bigskip


\subsection{오래된 비방에 대한 반박}


먼저 오래된 비난과 최초의 고발자들에게 답하고, 그 다음에 최근의 것에 답하겠습니다. 나를 고발한 자들은 오래전부터 많았고, 그들의 거짓 비난은 수년간 계속되었습니다. 나는 아뉘토스(Anytus)와 그 동료들보다 이 자들을 더 두려워합니다.


이 자들은 여러분이 어린아이였을 때부터 여러분의 마음을 장악했습니다. 그들은 이렇게 말했습니다: '소크라테스라는 자가 있는데, 하늘 위의 것과 땅 아래의 것을 탐구하고, 약한 논변을 강하게 만든다.' 이런 소문을 퍼뜨리는 자들이야말로 내가 두려워하는 고발자입니다. 이런 탐구를 하는 자는 신을 믿지 않는다고 사람들이 생각하기 쉽기 때문입니다.


그리고 이들의 이름을 나는 알 수도, 말할 수도 없습니다 — 희극 시인의 경우를 제외하면.


\tnote{'희극 시인'은 아리스토파네스(Aristophanes)를 가리킴. 그의 희극 〈구름(Νεφέλαι)〉에서 소크라테스를 공중에 매달려 헛소리하는 자연철학자로 희화화했음.}


\bigskip


\subsection{소크라테스의 지혜 — 델포이 신탁}


내가 교사라거나 돈을 받는다는 보고도 근거가 없습니다. 물론 누군가가 가르칠 수 있다면, 보수를 받는 것은 존경할 만한 일이라고 생각합니다. 레온티움의 고르기아스, 케오스의 프로디코스, 엘리스의 히피아스 같은 이들이 있지요.


나에 대한 이런 비난의 기원은 무엇인가? 공정한 질문이니 설명하겠습니다. 아테네 시민 여러분, 나의 이 평판은 내가 가진 일종의 지혜에서 비롯되었습니다. 어떤 종류의 지혜냐고요? 인간이 도달할 수 있는 정도의 지혜입니다. 그 정도까지는 내가 지혜롭다고 생각합니다. 반면 내가 방금 말한 사람들은 초인적인(superhuman) 지혜를 가졌을 수도 있는데, 나는 그것을 갖고 있지 않으므로 묘사할 수가 없습니다.


여기서 방해하지 말아주십시오. 내가 할 말이 과장되게 들릴 수 있지만, 이 말은 내 것이 아닙니다. 내 지혜에 대해 증언해줄 증인은 델포이의 신입니다.


\tnote{'델포이의 신'은 아폴론(Apollo). 델포이 신전의 신탁(oracle)은 아폴론의 여사제 퓌티아(Pythia)를 통해 전달됨.}


여러분은 카이레폰(Chaerephon)을 아실 것입니다. 그는 젊을 때부터 나의 친구였고, 여러분 대중의 친구이기도 했습니다. 카이레폰은 모든 일에 대단히 충동적인 사람이었는데, 대담하게도 델포이에 가서 신탁에게 나보다 더 지혜로운 자가 있느냐고 물었습니다. 퓌티아 여사제는 나보다 더 지혜로운 자는 없다고 답했습니다. 카이레폰 자신은 돌아가셨지만, 법정에 있는 그의 형제가 이 이야기의 진실을 확인해줄 것입니다.


왜 이것을 언급하느냐? 내가 왜 이런 나쁜 평판을 얻었는지 설명하기 위해서입니다. 이 답을 듣고 나는 자문했습니다: '신은 무엇을 뜻하는가? 이 수수께끼의 해석은 무엇인가?' 나는 크든 작든 어떤 지혜도 없음을 알고 있었기 때문입니다. 그가 내가 가장 지혜로운 사람이라고 할 때, 무엇을 뜻하는 것인가? 그는 신이니 거짓말을 할 수 없다 — 그것은 그의 본성에 어긋나니까.


오랜 숙고 끝에 이 문제를 시험할 방법을 생각해냈습니다. 나보다 지혜로운 사람을 찾으면, 그것을 증거로 들고 신에게 가면 되겠다고 생각했습니다.


\bigskip


\subsection{정치인, 시인, 장인을 찾아가다}


그래서 지혜롭다고 평판이 자자한 한 사람에게 갔습니다 — 정치인이었는데, 이름은 언급하지 않겠습니다. 결과는 이러했습니다: 그와 대화해보니, 많은 사람들이 그가 지혜롭다고 생각하고 그 자신은 더더욱 그렇게 생각했지만, 실은 지혜롭지 않다는 것이 느껴졌습니다. 그래서 그에게 스스로 지혜롭다고 생각하지만 실은 그렇지 않다는 것을 설명하려 했고, 그 결과 그는 나를 미워하게 되었습니다.


떠나면서 나는 속으로 생각했습니다: '우리 둘 다 정말로 아름답고 좋은 것은 아무것도 모르지만, 나는 그보다 나은 처지에 있다 — 그는 아무것도 모르면서 안다고 생각하지만, 나는 모르고 또한 안다고 생각하지도 않으니까. 이 점에서 나는 그보다 약간 유리한 것 같다.'


이후 사람들을 하나하나 찾아다녔습니다. 내가 원한을 산다는 것을 의식하지 못한 것은 아니며, 이를 걱정하고 두려워했습니다. 하지만 신의 말씀을 먼저 따라야 한다고 생각했습니다.


\tnote{'신의 말씀(the word of God)'은 Jowett 번역의 기독교적 색채. 원문은 'to tou theou'로, '신의 것/신에 관한 것'에 가까움.}


아테네 시민 여러분, 개의 이름을 걸고 맹세합니다! 진실을 말해야 하니까요 — 나의 탐구의 결과는 이것이었습니다: 가장 평판이 높은 자들이 거의 가장 어리석었고, 열등하다고 여겨지는 자들이 오히려 더 지혜롭고 나았습니다.


\tnote{'개의 이름을 걸고(by the dog)' — 그리스어 'nē ton kuna(νὴ τὸν κύνα)'. 소크라테스 특유의 맹세. 신의 이름을 직접 걸지 않으려는 우회적 표현으로 해석되기도 하고, 이집트 신 아누비스를 가리킨다는 설도 있음. 정확한 의미는 불확실.}


정치인들을 떠나 시인들에게 갔습니다 — 비극, 디튀람보스, 모든 종류의 시인들. 그들의 가장 정교한 구절을 가져가 그 의미를 물었습니다 — 그들이 나에게 무언가를 가르쳐줄 것이라 생각하면서. 거의 부끄러울 정도이지만 말해야 합니다: 이 자리에 있는 거의 모든 사람이 그들 자신보다 그들의 시에 대해 더 잘 말했을 것입니다. 이것이 보여준 것은, 시인들은 지혜에 의해서가 아니라 일종의 천부적 재능과 영감(genius and inspiration)에 의해 시를 쓴다는 것입니다. 그들은 점쟁이나 예언자와 같아서, 좋은 말을 많이 하지만 자기가 하는 말의 의미를 이해하지는 못합니다. 게다가 시인들은 시를 잘 쓴다는 이유로, 다른 모든 분야에서도 자신이 가장 지혜롭다고 믿었습니다.


마지막으로 장인(artisans)들에게 갔습니다. 나는 사실상 아무것도 모른다는 것을 의식하고 있었고, 그들은 좋은 것을 많이 알 것이라 확신했습니다. 이 점에서 나는 틀리지 않았습니다 — 그들은 내가 모르는 많은 것을 알고 있었고, 이 점에서 분명히 나보다 지혜로웠습니다. 하지만 훌륭한 장인들조차 시인들과 같은 오류에 빠져 있었습니다: 자기 기술(craft)에 뛰어나다는 이유로, 온갖 고상한 문제에 대해서도 안다고 생각한 것입니다. 이 결함이 그들의 지혜를 가렸습니다. 그래서 나는 신탁을 대신하여 자문했습니다: 그들의 지식도 무지도 없는 지금의 나로 있을 것인가, 아니면 그들처럼 둘 다 가진 채로 있을 것인가? 그리고 지금의 내가 더 나은 처지라고 답했습니다.


\bigskip


\subsection{신만이 지혜롭다}


이 탐구는 나에게 최악의, 가장 위험한 종류의 적을 많이 만들어주었고, 또한 많은 비방의 원인이 되었습니다. 나는 '지혜로운 자'로 불립니다 — 내 말을 듣는 사람들은 내가 남에게서 부족함을 발견하는 그 지혜를 내가 소유하고 있다고 상상하기 때문입니다.


\textbf{그러나 진실은, 아테네 시민 여러분, 신만이 지혜롭다는 것입니다(God only is wise). 그리고 이 신탁에서 신이 말하고자 한 것은, 인간의 지혜는 거의 또는 전혀 가치가 없다는 것입니다. 신은 소크라테스에 대해 말하는 것이 아닙니다. 다만 나의 이름을 예시(illustration)로 사용하고 있을 뿐입니다. 마치 이렇게 말하는 것처럼: '이 사람들 중에서 가장 지혜로운 자는, 소크라테스처럼, 자신의 지혜가 진실로(in truth) 아무런 가치가 없음을 아는 자이다.'}


\tnote{이 구절이 핵심. 원문의 'ho theos monos sophos'는 '신만이 sophos(지혜로운 자)이다'. Jowett의 'God only is wise'에서 'God'은 기독교적 유일신의 뉘앙스를 줄 수 있으나, 원문은 아폴론 또는 신적 존재 일반을 가리킬 수 있음. 'sophos'도 단순히 '똑똑한'이 아니라 '참된 앎을 가진' 상태를 뜻함.}


그래서 나는 신에게 순종하여 길을 갑니다. 시민이든 이방인이든 지혜로워 보이는 사람에게 다가가 탐구합니다. 그가 지혜롭지 않으면, 신탁을 입증하기 위해 그가 지혜롭지 않음을 보여줍니다. 이 작업에 완전히 몰두하여, 공적인 일에도 사적인 일에도 시간을 쓸 여유가 없고, 신에 대한 헌신 때문에 극심한 가난 속에 있습니다.


또 하나 있습니다 — 부유한 집안의 젊은이들이 할 일이 별로 없어서 자발적으로 나를 따라다닙니다. 그들은 허세부리는 자들이 검토당하는 것을 듣기 좋아하고, 종종 나를 흉내 내어 다른 사람들을 직접 검토합니다. 그러면 검토당한 자들은 자기 자신에게 화내는 대신 나에게 화를 냅니다: '이 저주스러운 소크라테스, 이 청년을 타락시키는 악당!' 하고요. 그리고 누가 '어떤 악행을 가르치느냐?'고 물으면, 그들은 모르고, 답할 수도 없습니다. 그래서 난처해 보이지 않기 위해 모든 철학자에게 갖다 붙이는 기성품 비난을 반복합니다: '구름 속의 것을 가르치고, 신을 믿지 않으며, 약한 논변을 강하게 만든다.'


이것이 나의 세 고발자 — 멜레토스, 아뉘토스, 뤼콘 — 이 나를 공격한 이유입니다. 멜레토스는 시인들을 대신하여, 아뉘토스는 장인들을 대신하여, 뤼콘은 수사학자들을 대신하여. 아테네 시민 여러분, 이것이 진실이며 온전한 진실입니다. 나는 아무것도 숨기지 않았고, 아무것도 위장하지 않았습니다. 이 솔직함이 그들로 하여금 나를 미워하게 만든다는 것을 알고 있지만, 그 미움 자체가 내가 진실을 말하고 있다는 증거 아닙니까?


\bigskip


\subsection{멜레토스에 대한 반박 — 청년 타락 혐의}


이제 멜레토스가 이끄는 두 번째 부류의 고발자들에게로 넘어갑니다. 그가 스스로 칭하기를 '선량하고 애국적인 사람'인 멜레토스. 그들의 고발장은 대략 이렇습니다: '소크라테스는 악행을 저지르는 자이며, 청년을 타락시키고, 국가의 신을 믿지 않으며, 자기만의 새로운 신적 존재를 갖고 있다.'


\noindent\textbf{소크라테스}: 이리 와서 질문에 답하시오, 멜레토스. 당신은 청년의 개선에 대해 많이 생각한다고?


\noindent\textbf{멜레토스}: 그렇습니다.


\noindent\textbf{소크라테스}: 그러면 재판관들에게 말하시오. 누가 청년을 개선하는 자인가? 그들을 타락시키는 자를 발견하는 수고를 했으니, 개선하는 자도 알 것 아닌가.


(멜레토스 침묵)


\noindent\textbf{소크라테스}: 보시오, 멜레토스, 당신은 침묵하고 할 말이 없소. 이것이야말로 당신이 이 문제에 관심이 없다는 상당한 증거가 아니겠소?


\noindent\textbf{멜레토스}: 법률이오.


\noindent\textbf{소크라테스}: 그건 내 질문의 뜻이 아니오. 법률을 아는 사람이 누구냐는 것이오.


\noindent\textbf{멜레토스}: 여기 재판정에 있는 재판관들이오, 소크라테스.


\noindent\textbf{소크라테스}: 이 재판관들이 청년을 교육하고 개선할 수 있다는 뜻이오?


\noindent\textbf{멜레토스}: 물론이오.


\noindent\textbf{소크라테스}: 전부 다? 아니면 일부만?


\noindent\textbf{멜레토스}: 전부 다.


\noindent\textbf{소크라테스}: 헤라 여신의 이름으로, 좋은 소식이로다! 개선하는 자가 참으로 많군. 그러면 청중은? 그들도 청년을 개선하는가?


\noindent\textbf{멜레토스}: 그렇소.


\noindent\textbf{소크라테스}: 원로원 의원들은?


\noindent\textbf{멜레토스}: 원로원 의원들도 개선하오.


\noindent\textbf{소크라테스}: 시민 의회 구성원들은 청년을 타락시키는가? 아니면 그들도 개선하는가?


\noindent\textbf{멜레토스}: 개선하오.


\noindent\textbf{소크라테스}: 그러면 모든 아테네인이 청년을 개선하고 향상시킨다 — 나 하나만 제외하고. 나 혼자만이 그들을 타락시킨다는 것인가? 당신이 주장하는 것이 그것인가?


\noindent\textbf{멜레토스}: 그것을 강력히 주장하오.


소크라테스는 답합니다: 그것이 사실이라면 나는 참으로 불행한 자이오. 하지만 질문 하나 하겠소: 말(horse)의 경우에도 그러한가? 한 사람이 해를 끼치고 온 세상이 이로움을 주는가? 정확히 그 반대 아닌가? 한 사람 — 말 조련사 — 이 이로움을 줄 수 있고, 나머지는 오히려 해를 끼치지 않는가?


소크라테스는 이어서 묻습니다: 좋은 이웃은 이웃에게 이로움을 주고, 나쁜 이웃은 해를 끼친다는 것을 인정하지 않는가? 그렇다면, 내가 함께 사는 사람을 타락시키면 그에 의해 해를 입을 가능성이 높다는 것을 — 70세가 넘은 내가 모를 리가 있겠는가? 그런데 내가 고의로 그런다고? 당신의 논리대로라면 나는 둘 중 하나이오: 타락시키지 않거나, 비의도적으로 타락시키거나. 비의도적이라면 법은 그것을 다루지 않소 — 당신은 나를 사적으로 불러 충고했어야 했소. 법정은 교육의 장소가 아니라 처벌의 장소이오.


\bigskip


\subsection{멜레토스에 대한 반박 — 무신론 혐의}


\noindent\textbf{소크라테스}: 멜레토스, 좀 더 분명히 말해보시오. 내가 어떤 신들을 가르치되 국가의 신과 다른 신이라는 것인가, 아니면 내가 완전한 무신론자라는 것인가?


\noindent\textbf{멜레토스}: 후자요 — 당신은 완전한 무신론자요.


\noindent\textbf{소크라테스}: 그것은 놀라운 주장이오. 내가 태양이나 달의 신성을 믿지 않는다는 뜻인가? 그것은 모든 사람의 공통된 믿음인데?


\noindent\textbf{멜레토스}: 재판관 여러분, 그는 믿지 않습니다! 그는 태양은 돌이고 달은 흙이라고 말합니다!


\noindent\textbf{소크라테스}: 멜레토스여, 당신은 클라조메나이의 아낙사고라스를 고발하고 있는 것이오! 재판관들이 그 교리가 아낙사고라스의 책에 있다는 것을 모를 만큼 무지하다고 생각하는 것이오? 이 교리는 극장에서 입장료 1 드라크마에 종종 전시되는데, 그것을 내가 만들어냈다고 하면 청년들이 나를 비웃을 것이오.


\tnote{아낙사고라스(Anaxagoras)는 소크라테스 이전 세대의 자연철학자로, 태양은 불타는 돌이라고 주장하여 불경죄로 아테네에서 추방당했음.}


소크라테스는 이어서 논리적 함정을 드러냅니다: 고발장에서 멜레토스는 소크라테스가 신적/영적 존재(divine or spiritual agencies)를 가르치고 믿는다고 맹세했습니다. 그런데 영적 존재를 믿으면서 신을 믿지 않을 수 있는가? 영적 존재(daimones)는 신이거나 신의 자식이 아닌가? 노새의 존재는 인정하면서 말과 당나귀의 존재를 부정하는 것과 같은 말도 안 되는 소리이오.


\tnote{'daimones(다이모네스)'는 단순히 '악마'가 아니라 신적 존재/영적 존재 전반을 가리킴. Jowett의 'demigods'도 정확하지 않음.}


\bigskip


\subsection{죽음에 대하여 — 철학자의 사명}


누군가 물을 것입니다: '소크라테스, 당신의 생활방식이 때 이른 죽음을 가져올 수 있는데 부끄럽지 않은가?' 이에 답하겠습니다: 당신이 틀렸소. 무엇이든 쓸모 있는 사람은 살거나 죽을 확률을 계산해서는 안 되오. 자신이 옳은 일을 하고 있는지 그른 일을 하고 있는지 — 좋은 사람의 역할을 하는지 나쁜 사람의 역할을 하는지 — 만 생각해야 하오.


\textbf{죽음에 대한 두려움이야말로 지혜의 가장(pretence)이지 실제 지혜가 아닙니다 — 알지 못하는 것을 아는 것처럼 꾸미는 것입니다. 아무도 죽음이 최대의 악인지 모르니까요 — 어쩌면 최대의 선일 수도 있습니다. 이런 무지의 허세야말로 부끄러운 종류의 무지 아닙니까? 그리고 이 점에서, 내가 다른 사람들보다 나을 수 있다고 생각하는 점이 있다면, 저승에 대해 거의 모르면서도 안다고 생각하지 않는다는 것입니다. 다만 불의와, 신이든 인간이든 자신보다 나은 존재에 대한 불복종이 악이며 불명예라는 것은 알고 있습니다.}


\tnote{'pretence of wisdom' — 원문에서 소크라테스는 죽음을 두려워하는 것 자체를 '모르는 것을 아는 척하는 것'이라 규정함. 이것이 그의 인식론적 겸손의 핵심 구조.}


그래서 만약 여러분이 나를 풀어주되, 더 이상 탐구하고 사색하지 말라는 조건을 붙이고, 다시 걸리면 죽일 것이라 한다면 — 나는 이렇게 답하겠습니다: 아테네 시민 여러분, 나는 여러분을 존경하고 사랑합니다. 하지만 나는 여러분보다 신에게 복종하겠습니다. 살아있고 힘이 있는 한, 나는 철학의 실천과 가르침을 결코 그만두지 않을 것입니다.


\tnote{'나는 여러분보다 신에게 복종하겠습니다(I shall obey God rather than you)' — 원문 'peissomai tōi theōi mallon ē hymin'. 이 선언은 후대 기독교에서도 자주 인용되었으나, 원래 맥락은 아테네의 다신교적 세계관 속에서의 발언.}


나는 만나는 모든 사람에게 이렇게 말합니다: '나의 친구여, 위대하고 강하며 지혜로운 도시 아테네의 시민인 당신이, 돈과 명예와 평판을 최대한 쌓는 것에는 그토록 관심을 쏟으면서, 지혜와 진리와 영혼의 최대한의 개선에는 왜 이토록 무관심한가? 부끄럽지 않은가?'


이것이 신의 명령이니, 여러분도 아시기 바랍니다. 오늘까지 신에 대한 나의 봉사보다 더 큰 선이 이 나라에 일어난 적은 없다고 믿습니다. 나는 여러분 모두를, 젊은이든 늙은이든, 설득하러 다닐 뿐입니다 — 먼저 자신의 몸이나 재산이 아니라, 영혼의 최대한의 개선을 돌보라고. 덕(virtue)은 돈에서 나오지 않으며, 덕에서 돈과 모든 다른 좋은 것이 나옵니다. 이것이 나의 가르침이고, 만약 이것이 청년을 타락시키는 교리라면, 나의 영향은 실로 파멸적이겠지요.


\textbf{아테네 시민 여러분, 아뉘토스가 권하는 대로 하든 말든, 무죄로 하든 유죄로 하든, 알아두십시오: 나는 결코 나의 방식을 바꾸지 않을 것입니다. 여러 번 죽어야 하더라도.}


\bigskip


\subsection{등에/쇠파리(gadfly) 비유}


하나 더 말씀드리겠습니다. 여러분이 나 같은 사람을 죽이면, 나보다 여러분 자신이 더 큰 해를 입을 것입니다. 나를 죽이면, 여러분은 나 같은 사람을 쉽게 찾지 못할 것입니다. 우스꽝스러운 비유를 쓸 수 있다면, 나는 일종의 등에(gadfly, 쇠파리)입니다 — 신이 이 나라에 준 것입니다. 이 나라는 크고 고귀한 말과 같은데, 그 크기 때문에 움직임이 굼뜨고, 삶으로 자극받을 필요가 있습니다. 내가 바로 그 등에이며, 신이 이 나라에 붙여놓은 것입니다. 하루 종일, 어디서나, 항상 여러분에게 달라붙어 깨우고, 설득하고, 질책합니다.


내가 신으로부터 여러분에게 주어진 존재임은 이것으로 증명됩니다: 내가 다른 사람들과 같았다면, 이 모든 세월 동안 내 일을 소홀히 하면서 여러분의 일을 돌보고, 아버지나 형처럼 한 사람 한 사람에게 다가가 덕을 돌보라고 권하지는 않았을 것입니다 — 이것은 인간의 본성에 맞지 않습니다. 그리고 내가 이것으로 이득을 얻었다면 그래도 이해가 가겠지만, 고발자들조차 내가 보수를 받거나 구한 적이 있다고 감히 말하지 못합니다. 내 가난이 이 진실의 충분한 증인입니다.


\bigskip


\subsection{왜 공적인 삶에 나서지 않았는가}


누군가 궁금할 것입니다: 사적으로 돌아다니며 남의 일에 관여하면서, 왜 공개적으로 나서서 국가에 조언하지 않는가? 그 이유를 말씀드리겠습니다. 여러분은 내가 종종 말하는 것을 들었습니다 — 나에게 오는 신탁 또는 표징(sign)이 있다고. 이것은 멜레토스가 고발장에서 조롱하는 바로 그 신적인 것(divinity)입니다. 이 표징은 어릴 때부터 있었습니다. 이 표징은 내가 무언가를 하려고 할 때 금하는 목소리인데, 무언가를 하라고 명령하지는 않습니다.


\tnote{'다이모니온(daimonion)' — 소크라테스의 내면의 목소리. '개인적 신적 표징'으로, 제도화된 종교와는 다른 것. 이것이 '새로운 신적 존재를 도입한다'는 혐의의 근거가 됨. 다이모니온의 정확한 성격(양심의 소리, 종교적 체험, 수사적 장치 등)에 대해서는 학자마다 견해가 다름.}


내가 정치에 참여했다면 오래전에 죽었을 것이고, 여러분에게도 나 자신에게도 아무런 좋은 일을 하지 못했을 것입니다. 진실을 말하겠습니다: 대중과 전쟁하면서 정의를 위해 정직하게 싸우는 사람은 목숨을 부지할 수 없습니다. 잠시라도 살고자 하는 사람은 사적인 위치에 머물러야지, 공적인 위치에 있어서는 안 됩니다.


\bigskip


\subsection{실천의 증거 — 두 가지 사건}


말이 아니라 행동으로 증거를 보여드리겠습니다.


내가 국가에서 유일하게 맡았던 공직은 원로원 의원이었습니다. 내 부족인 안티오키스 부족이 아르기누사이 해전 이후 전사자의 시신을 수습하지 못한 장군들의 재판에서 의장직을 맡았을 때, 여러분은 그들을 한꺼번에 재판하자고 했습니다 — 이것은 불법이었고, 여러분 모두도 나중에 그렇게 생각했습니다. 그때 프뤼타네이스 중에서 나 혼자만이 이 불법에 반대했고, 반대표를 던졌습니다. 연설가들이 나를 탄핵하고 체포하겠다고 위협했고, 여러분이 소리치고 외쳤지만, 나는 감옥과 죽음이 두려워 여러분의 불의에 가담하기보다는 법과 정의 편에서 위험을 감수하기로 했습니다.


그리고 30인 과두정이 권력을 잡았을 때, 그들은 나와 네 명을 원형홀(rotunda)로 소환하여 살라미스 사람 레온을 살라미스에서 데려오라고 명령했습니다 — 그를 처형하려고. 이런 종류의 명령은 가능한 한 많은 사람을 범죄에 연루시키려는 의도였습니다. 그때 나는 말이 아니라 행동으로 보여주었습니다 — 죽음에 대해서는 전혀 개의치 않으며, 나의 유일한 두려움은 불의하거나 불경한 일을 하는 것이라고. 그 억압적 권력의 강한 팔도 나를 두렵게 하여 불의를 행하게 만들지 못했습니다. 원형홀을 나왔을 때 나머지 네 명은 살라미스로 가서 레온을 데려왔지만, 나는 조용히 집으로 갔습니다. 30인 정권이 곧 무너지지 않았다면, 나는 목숨을 잃었을 것입니다.


\bigskip


\subsection{소크라테스의 제자들에 대하여}


나는 모든 행동에서, 공적으로나 사적으로나, 항상 같은 사람이었습니다. 나에게는 정규 제자(disciples)가 없습니다. 다만 누구든 나의 사명 수행을 듣고 싶으면, 젊든 늙든, 자유롭게 올 수 있습니다. 돈을 내는 사람하고만 대화하고 내지 않는 사람은 거부하는 것도 아닙니다. 누구든, 부자든 가난한 자든, 질문하고 답하고 내 말을 들을 수 있습니다.


만약 내가 정말로 청년을 타락시켰다면, 성장하여 제정신이 된 자들이 고발자로 나서서 복수해야 하지 않겠습니까? 그들이 직접 나오지 않더라도, 그들의 친족 — 아버지, 형제 — 이 나서야 하지 않겠습니까? 하지만 보십시오 — 이 자리에 많은 사람들이 있습니다.


크리톤(Crito)이 있고, 그의 아들 크리토불로스(Critobulus)도 보입니다. 스펫토스의 뤼사니아스(Lysanias)는 아이스키네스(Aeschines)의 아버지입니다. 케피소스의 안티폰(Antiphon)은 에피게네스(Epigenes)의 아버지입니다. 아리스톤(Ariston)의 아들 아데이만토스(Adeimantus)도 있는데, 그의 형제 플라톤(Plato)이 여기 있습니다. 아폴로도로스(Apollodorus)의 형제인 아이안토도로스(Aeantodorus)도 보입니다.


\tnote{플라톤이 재판에 참석했음을 소크라테스 본인의 입을 통해 확인하는 대목.}


이들 모두가 '타락시키는 자'인 나를 위해 증언할 준비가 되어 있습니다. 타락당한 청년들에게는 동기가 있을 수 있겠지만, 타락당하지 않은 연장자 친족들은 진실과 정의를 위해서가 아니면 왜 나를 지지하겠습니까?


\bigskip


\subsection{눈물과 탄원에 대하여}


아마 여러분 중 누군가는 불쾌할 것입니다 — 자신은 덜 심각한 경우에도 눈물과 탄원에 호소하고, 아이들을 법정에 데려와 동정을 구했는데, 나는 목숨이 위태로우면서도 이런 짓을 하지 않으니까.


나는 사람이고, 다른 사람들처럼 살과 피의 존재입니다 — 호메로스가 말하듯 나무나 돌이 아닙니다. 나에게도 가족이 있고, 아들이 셋입니다 — 하나는 성장 중이고 둘은 아직 어립니다. 하지만 무죄를 구걸하기 위해 그들을 이리 데려오지는 않겠습니다.


나이가 들어 지혜롭다는 평판이 있는 사람이 — 그 평판이 마땅하든 아니든 — 스스로를 비하해서는 안 됩니다. 그리고 명예의 문제를 떠나서, 재판관에게 탄원하여 무죄를 얻어내는 것은 잘못입니다. 재판관의 의무는 정의를 선물로 주는 것이 아니라, 판결을 내리는 것이기 때문입니다.


아테네 시민 여러분, 나의 주장과 여러분에게 가장 좋은 것이 무엇인지를 신에게 맡기겠습니다.


\begin{center}\textit{(배심원단이 소크라테스에게 유죄를 선고하다.)}\end{center}


\bigskip


\newpage
\section{제2부: 형량에 대한 소크라테스의 제안}


아테네 시민 여러분, 유죄 평결에 대해 슬퍼하지 않는 이유가 많습니다. 이미 예상했고, 다만 표차가 이렇게 근소하리라고는 생각하지 못했습니다. 30표만 다른 쪽으로 넘어갔더라면 무죄가 되었을 것입니다.


멜레토스는 사형을 형벌로 제안합니다. 그러면 나는 나의 몫으로 무엇을 제안해야 하겠습니까? 분명히 내가 마땅히 받아야 할 것을. 평생 한가히 있을 지혜가 없었던 사람 — 대중이 관심 갖는 것, 재산, 가문의 이익, 군직, 의회 연설, 관직, 음모, 당파를 돌보지 않은 사람에게 무엇을 해야 하겠습니까? 나는 너무 정직한 사람이어서 이런 길을 따라가면 살아남을 수 없다고 판단했고, 여러분 한 사람 한 사람에게 가장 큰 선을 사적으로 행할 수 있는 곳으로 갔습니다 — 각자 자기 자신을 돌보고, 사적 이익보다 덕과 지혜를 먼저 구하라고 설득했습니다.


\textbf{이런 사람에게 무엇을 해야 합니까? 의심할 것 없이 좋은 일이겠지요. 가난하지만 여러분의 은인이며, 여러분을 가르칠 수 있도록 여가가 필요한 사람에게 합당한 보상은 무엇이겠습니까? 프뤼타네이온(Prytaneum)에서의 국비 식사보다 더 적절한 보상은 없습니다. 올림피아에서 마차 경주로 우승한 시민보다 훨씬 더 이 보상을 받을 자격이 있습니다. 그는 여러분에게 행복의 외관만을 주지만, 나는 실체를 줍니다.}


\tnote{프뤼타네이온은 아테네의 공공건물로, 올림피아 우승자 등 국가 공로자에게 국비 식사가 제공되던 곳. 소크라테스가 사형 대신 이것을 제안한 것은 도발적 행위.}


사형이 선인지 악인지 모르는데, 확실히 악인 형벌을 왜 제안해야 하겠습니까? 감옥을 말해야 합니까? 벌금을? 돈이 없는데 같은 이야기입니다. 추방을? 내 동료 시민인 여러분이 나의 담론을 견디지 못한다면, 다른 곳 사람들이 견딜 것 같습니까?


누군가 말할 것입니다: '소크라테스, 입을 다물고 있을 수는 없는가?' 이것이야말로 여러분이 이해하기 가장 어려운 나의 답입니다. 입을 다무는 것은 신의 명령에 대한 불복종이라고 말하면, 여러분은 내가 진지하다고 믿지 않을 것입니다. \textbf{그리고 인간에게 가장 큰 선은 매일 덕에 대해 대화하는 것이며, 검토되지 않은 삶은 살 가치가 없다고 말하면 — 여러분은 더더욱 믿지 않을 것입니다.}


\tnote{'검토되지 않은 삶은 살 가치가 없다(the unexamined life is not worth living)' — 원문 'ho de anexetastos bios ou biōtos anthrōpōi'. 서양 철학에서 가장 많이 인용되는 문장 중 하나.}


하지만 내가 하는 말은 사실이며, 여러분을 설득하기 어려운 것이기도 합니다. 1 므나(mina) 정도는 낼 수 있으니 그것을 제안합니다. 여기 있는 플라톤, 크리톤, 크리토불로스, 아폴로도로스가 30 므나를 내겠다고 합니다. 그러면 30 므나를 형벌로 제안합니다.


\begin{center}\textit{(배심원단이 소크라테스에게 사형을 선고하다.)}\end{center}


\bigskip


\newpage
\section{제3부: 사형 선고 후 소크라테스의 말}


\subsection{사형에 투표한 자들에게}


아테네 시민 여러분, 나를 죽이는 대가로 얻는 시간은 얼마 되지 않을 것이고, 여러분은 도시의 비방자들로부터 '지혜로운 소크라테스를 죽였다'는 악명을 얻게 될 것입니다 — 비록 내가 지혜롭지 않더라도, 여러분을 비난하고 싶을 때는 그렇게 부를 것입니다. 조금만 기다렸더라면, 자연의 과정 속에서 여러분의 바람은 이루어졌을 것입니다. 보시다시피 나는 나이가 많이 들었고 죽음이 멀지 않습니다.


여러분은 내가 말의 부족으로 유죄판결을 받았다고 생각합니다. 아닙니다. 부족한 것은 말이 아니었습니다. 부족한 것은 대담함, 뻔뻔함, 여러분이 좋아할 방식으로 말하려는 의지였습니다 — 울고, 탄식하고, 통곡하며, 다른 사람들에게서 듣는 데 익숙한 많은 말과 행동을 하는 것. 하지만 위험의 순간에 비천하거나 천박한 행동을 해서는 안 된다고 생각했습니다.


\textbf{나의 방식대로 말하고 죽는 것이, 여러분의 방식대로 말하고 사는 것보다 낫습니다.}


전쟁이든 법정이든, 죽음을 피하기 위해 모든 수단을 동원해서는 안 됩니다. 어려운 것은 죽음을 피하는 것이 아니라, 불의(unrighteousness)를 피하는 것입니다 — 불의는 죽음보다 빨리 달리니까요. 나는 늙고 느려서 느린 주자에게 따라잡혔고, 나의 고발자들은 날카롭고 빠르지만 더 빠른 주자인 불의에 따라잡혔습니다.


그리고 이제 나는 여러분에 의해 사형을 선고받고 떠나며, 그들도 진실에 의해 악덕과 불의의 형벌을 선고받고 떠납니다.


\bigskip


\subsection{예언}


그리고 이제, 나를 사형에 처한 사람들에게 예언하겠습니다. 죽음이 다가오면 사람에게 예언의 능력이 주어지는 법이니까요. 나를 살해한 자들이여, 내가 죽은 직후 여러분이 나에게 가한 것보다 훨씬 무거운 벌이 여러분을 기다리고 있다고 예언합니다. 여러분은 고발자를 피하고, 자신의 삶에 대한 해명을 하지 않으려고 나를 죽였습니다. 하지만 결과는 여러분이 생각하는 것과 정반대일 것입니다. 지금까지 내가 억제해 온 고발자들이 더 많이 나타날 것이고, 더 젊으니 더 가혹할 것이며, 여러분은 더 분노하게 될 것입니다.


\textbf{사람을 죽여서 삶을 비판하는 자를 피할 수 있다고 생각한다면 틀렸습니다. 그런 탈출은 가능하지도 명예롭지도 않습니다. 가장 쉽고 고귀한 길은 다른 사람을 짓누르는 것이 아니라 자기 자신을 개선하는 것입니다.}


\bigskip


\subsection{무죄에 투표한 친구들에게}


나의 무죄에 투표해주었을 친구들이여, 내게 일어난 이 일의 의미에 대해 이야기하고 싶습니다. 지금까지 내 안의 익숙한 신탁 — 다이모니온 — 은 사소한 일에도 내가 실수하려 하면 반대하곤 했습니다. 그런데 보시다시피, 일반적으로 최후의 최악의 악으로 여겨지는 이 일이 닥쳤는데, 신탁은 아침에 집을 나설 때도, 이 법정에 올라올 때도, 말하는 동안에도, 어떤 반대의 표징도 보이지 않았습니다.


나는 이것을 이렇게 해석합니다: 나에게 일어난 일은 선(good)이며, 죽음이 악이라고 생각하는 우리가 틀렸다는 것입니다.


\bigskip


\subsection{죽음에 대한 희망}


다른 방식으로도 생각해 봅시다. 죽음이 선이라고 희망할 큰 이유가 있습니다. 둘 중 하나이니까요: 죽음은 허무와 완전한 무의식의 상태이거나, 사람들이 말하듯 영혼이 이 세계에서 다른 세계로 이주하는 것입니다.


의식이 없는 상태라면 — 꿈에 의해서도 방해받지 않는 잠과 같다면 — 죽음은 헤아릴 수 없는 이득일 것입니다. 꿈도 없이 잔 밤과 삶의 다른 낮과 밤을 비교하여, 그보다 더 좋고 유쾌하게 보낸 날밤이 몇이나 되는지 세어보면 — 보통 사람은 말할 것도 없고 대왕(the great king)조차도 그런 날밤을 많이 찾지 못할 것입니다. 죽음이 이런 것이라면, 죽는 것은 이득입니다 — 영원이 단 하나의 밤에 불과하니까요.


\tnote{'대왕(the great king)'은 페르시아 왕을 가리킴. 당시 아테네인에게 가장 부유하고 권력 있는 존재의 상징.}


하지만 죽음이 다른 곳으로의 여행이라면, 그곳에서 모든 죽은 자를 만날 수 있다면 — 이보다 더 큰 선이 어디 있겠습니까? 미노스, 라다만튀스, 아이아코스, 트립톨레모스, 그리고 생전에 의로웠던 다른 신의 아들들이 심판하는 곳에 도착한다면, 그 순례는 할 가치가 있을 것입니다. 오르페우스, 무사이오스, 헤시오도스, 호메로스와 대화할 수 있다면 무엇을 내놓지 않겠습니까?


무엇보다, 저 세계에서도 이 세계에서처럼 참된 지식과 거짓 지식에 대한 탐구를 계속할 수 있을 것입니다. 누가 지혜롭고 누가 지혜로운 척만 하는지 알아낼 수 있을 것입니다. 저 세계에서는 이것 때문에 사람을 죽이지 않으니까요.


\bigskip


\subsection{마지막 말}


그러므로, 재판관들이여, 죽음에 대해 낙관하시고, 이 진실을 아십시오 — 좋은 사람에게는 살아서나 죽어서나 어떤 악도 일어날 수 없습니다. 그와 그의 것은 신들에게 방치되지 않습니다. 나의 다가오는 최후도 단순한 우연이 아닙니다. 죽어서 풀려나는 것이 나에게 더 나았다는 것을 분명히 봅니다. 그래서 신탁도 아무런 표징을 주지 않았던 것입니다. 이런 이유로, 나는 고발자들에게도 유죄를 선고한 자들에게도 화내지 않습니다 — 그들은 나에게 아무런 해도 끼치지 않았습니다.


마지막으로 한 가지 부탁이 있습니다. 내 아들들이 자라면, 내가 여러분을 괴롭혔듯이 그들을 괴롭혀 주십시오 — 덕보다 재산이나 다른 것에 관심을 보이거나, 아무것도 아니면서 무언가인 척한다면, 내가 여러분을 질책했듯이 그들을 질책해 주십시오. 그렇게 해주시면, 나와 내 아들들은 여러분의 손에서 정의를 받은 것입니다.


\textbf{떠날 시간이 되었습니다. 우리는 각자의 길을 갑니다 — 나는 죽으러, 여러분은 살러. 어느 쪽이 더 나은지는 신만이 아십니다.}


\tnote{마지막 문장 원문: 'hopoteron de hēmōn erchetai epi ameinon pragma, adēlon panti plēn ē tōi theōi.' 직역하면 '우리 중 누가 더 나은 일로 향하는지는, 신 외에는 누구에게도 불분명하다.' Jowett의 'God only knows'는 의미상 정확하나 기독교적 색채가 있음.}


\bigskip


\begin{center}\textit{끝}\end{center}


\end{document}
